% Описание алгоритма corridor_sweeping (Anatoliy/swarm_mavic/corridor_sweeping/main.py)
% Компиляция: pdflatex corridor_sweeping.tex (рекомендуется пакет babel для русского)
\documentclass[12pt,a4paper]{article}

\usepackage[utf8]{inputenc}
\usepackage[T2A]{fontenc}
\usepackage[russian]{babel}
\usepackage{amsmath,amssymb,amsthm}
\usepackage{geometry}
\geometry{margin=2.5cm}
\usepackage{hyperref}
\hypersetup{colorlinks=true,linkcolor=blue,urlcolor=blue}

\title{Алгоритм \texttt{corridor\_sweeping}:\\
симуляция прохода роя по коридору между двумя движущимися якорями}
\author{Описание по коду \texttt{main.py} и \texttt{point\_mass.py}}
\date{}

\begin{document}
\maketitle

\begin{abstract}
В файле \texttt{corridor\_sweeping/main.py} реализована двумерная (по плоскости $Oxy$) модель роя материальных точек, взаимодействующих только с агентами и двумя \emph{якорями}, попадающими в радиус видимости. Управление задаётся релейно (знак ошибки) с нелинейной функцией $\tanh$, зависящей от расстояний до ближайших соседей по четырём направлениям в относительных координатах. Якоря задают границы коридора и движутся по заданным траекториям $p_L(t)$, $p_R(t)$. Ниже приведены математическая постановка, обозначения с пояснениями при первом появлении и логика шагов симуляции.
\end{abstract}

\tableofcontents
\newpage

%---------------------------------------------------------------------
\section{Назначение и общая схема}

\paragraph{Цель сценария.}
Смоделировать поведение $N$ агентов вдоль коридора, образованного двумя параллельно движущимися опорными точками (левый и правый \emph{якорь}). Агенты избегают столкновений и стремятся выстроиться, используя только локальную информацию о соседях в круге радиуса видимости.

\paragraph{Объекты модели.}
\begin{itemize}
  \item \textbf{Агент $i$} --- материальная точка с положением $p_i(t) \in \mathbb{R}^3$, скоростью $v_i(t) \in \mathbb{R}^3$ и приложенной силой $u_i(t) \in \mathbb{R}^3$. В коде сила фактически задаёт ускорение при массе $m_i = 1$ (см.\ раздел~\ref{sec:dynamics}).
  \item \textbf{Якорь} --- неуправляемая опорная точка $a_k(t)$, $k \in \{L,R\}$, движущаяся по заданной функции времени; для агентов якорь ведёт себя как особый сосед в зоне видимости.
\end{itemize}

\paragraph{Параметры сценария (глобальные константы и аргументы).}
При \textbf{первом упоминании} ниже переменные сопровождаются пояснением.

\begin{itemize}
  \item $N$ --- \texttt{n\_agents}: число агентов в рое.
  \item $T$ --- число итераций симуляции \texttt{n\_iters}; дискретное время $n = 0,1,\ldots,T-1$.
  \item $\Delta t$ --- \texttt{dt}: шаг интегрирования по времени между обновлениями состояния.
  \item $R_{\mathrm{vis}}$ --- \texttt{R\_vis}: радиус видимости; агент «видит» другого агента или якорь, если евклидово расстояние между ними не превышает $R_{\mathrm{vis}}$.
  \item $u_{\max}$ --- \texttt{u\_max}: модуль горизонтальной составляющей управления по каждой оси $Ox$ и $Oy$ (релейный закон даёт $\pm u_{\max}$).
  \item $w$ --- \texttt{w}: положительная константа, подставляемая в функцию $g$ для оси $Oy$ при отсутствии соседа в соответствующем направлении (см.\ \eqref{eq:g}); по оси $Ox$ в коде вместо неё используется $0$.
  \item $m_i$ --- \texttt{mass} класса \texttt{PointMass}: масса агента (по умолчанию $m_i = 1$).
\end{itemize}

\paragraph{Поток выполнения (\texttt{process}).}
На каждом шаге $n$:
\begin{enumerate}
  \item записать в лог состояния агентов и якорей;
  \item обновить время и положения якорей: $a_L(t+\Delta t)$, $a_R(t+\Delta t)$;
  \item для каждого живого агента $i$: найти множество соседей, вычислить управление $u_i$, применить силу и выполнить шаг интегрирования динамики на $\Delta t$.
\end{enumerate}

%---------------------------------------------------------------------
\section{Кинематика якорей}

Якорь задаётся функцией движения $f_k: \mathbb{R}_{\ge 0} \to \mathbb{R}^3$, $k \in \{L,R\}$ (в коде \texttt{left\_anchor\_motion\_func}, \texttt{right\_anchor\_motion\_func}).

\paragraph{Переменные.}
\begin{itemize}
  \item $t$ --- внутреннее время якоря; на шаге симуляции увеличивается: $t \leftarrow t + \Delta t$ (\texttt{AnchorWithFunc.step}).
  \item $a_k(t) = f_k(t)$ --- положение $k$-го якоря (\texttt{anchor.pose}).
\end{itemize}

Типичный сценарий в конце \texttt{main.py}: $f_L(t) = (v_x t,\, y_L,\, 0)^\top$, $f_R(t) = (v_x t,\, y_R,\, 0)^\top$ с одинаковой скоростью по $x$ и фиксированными $y_L$, $y_R$, то есть коридор параллелен оси $Ox$ и весь ансамбль переносится вдоль $x$.

%---------------------------------------------------------------------
\section{Инициализация агентов}

При создании агента $i$ его координата $z$ совпадает с $z$ правого якоря в начальный момент. Координата $y$ выбирается равномерно между $y$-координатами правого и левого якоря; координата $x$ --- равномерно на отрезке, выраженном через положения якорей и $R_{\mathrm{vis}}$ (см.\ исходный код \texttt{process}).

\paragraph{Переменные.}
\begin{itemize}
  \item $p_i^0 = (x_i^0, y_i^0, z_i^0)^\top$ --- начальное положение $i$-го агента.
  \item \texttt{id} --- целочисленный идентификатор агента для логирования и визуализации.
\end{itemize}

%---------------------------------------------------------------------
\section{Множество соседей и состояние \texttt{peer\_state}}
\label{sec:peers}

Для фокусного агента $i$ формируется список \emph{относительных векторов} к видимым объектам:
\begin{equation}
  \mathcal{P}_i = \{\, r_{ij} = p_j - p_i \;:\; \|r_{ij}\| \le R_{\mathrm{vis}},\; j \ne i,\; j \text{ жив} \,\}
  \;\cup\;
  \{\, r_{ik} = a_k - p_i \;:\; \|r_{ik}\| \le R_{\mathrm{vis}} \,\}.
\end{equation}

\paragraph{Переменные.}
\begin{itemize}
  \item $r = (r_x, r_y, r_z)^\top$ --- элемент списка \texttt{peers} (в коде \texttt{vec}); компоненты относительного смещения соседа к агенту $i$.
  \item $s_i \in \{0,1,2\}$ --- \texttt{peer\_state} агента $i$ после вызова \texttt{get\_peers}:
  \begin{itemize}
    \item $s_i = 0$, если нет видимых агентов и якорей (кроме самого себя);
    \item $s_i = 1$, если есть хотя бы один видимый агент;
    \item $s_i = 2$, если в зоне видимости есть хотя бы один якорь.
  \end{itemize}
  При наличии нескольких соседей состояние дополнительно обновляется как максимум по видимым агентам: $s_i \leftarrow \max(s_i, s_j)$ для всех видимых агентов $j$ (в коде \texttt{peer\_state = max(peer\_state, agent.peer\_state)}).
\end{itemize}

%---------------------------------------------------------------------
\section{Ближайшие соседи по четырём полупространствам}
\label{sec:nearest}

Для каждого $r = (r_x, r_y, r_z) \in \mathcal{P}_i$ вычисляются четыре «расстояния до ближайшего соседа» в смысле \emph{минимальной положительной компоненты} по осям:

\paragraph{Ось $x$.}
\begin{align}
  d_x^+ &= \min \{\, r_x \;:\; r \in \mathcal{P}_i,\; r_x \ge 0 \,\}, \\
  d_x^- &= \min \{\, -r_x \;:\; r \in \mathcal{P}_i,\; r_x < 0 \,\}
         = \min \{\, |r_x| \;:\; r_x < 0 \,\}.
\end{align}
Если множество пусто, в коде соответствующее значение остаётся бесконечным (\texttt{np.inf}), а флаг «сосед найден» --- ложным.

\paragraph{Ось $y$ (аналогично).}
\begin{align}
  d_y^+ &= \min \{\, r_y \;:\; r \in \mathcal{P}_i,\; r_y \ge 0 \,\}, \\
  d_y^- &= \min \{\, -r_y \;:\; r \in \mathcal{P}_i,\; r_y < 0 \,\}.
\end{align}

\paragraph{Переменные.}
\begin{itemize}
  \item $d_x^+$, $d_y^+$, $d_x^-$, $d_y^-$ --- скаляры, возвращаемые как \texttt{distances} (при отсутствии соседа --- $\infty$ в реализации).
  \item \texttt{flags} --- четыре булевых признака наличия соседа в каждом из четырёх направлений (\texttt{has\_peer\_x\_plus} и т.\,д.).
\end{itemize}

\textbf{Важно.} Величины $d_\cdot^\pm$ --- это не евклидовы расстояния до соседа, а минимальные проекции со знаком (для «$+$'' --- самая маленькая положительная компонента среди соседей в соответствующем полупространстве).

%---------------------------------------------------------------------
\section{Закон управления}
\label{sec:control}

\subsection{Вспомогательные функции}

\paragraph{Функция знака.}
\begin{equation}
  \operatorname{sgn}(\sigma) =
  \begin{cases}
    +1, & \sigma \ge 0, \\
    -1, & \sigma < 0.
  \end{cases}
\end{equation}
(В коде \texttt{sign}; нуль относится к положительной ветви.)

\paragraph{Нелинейность $\xi$.}
\begin{equation}
  \xi(g) = \tanh(g).
\end{equation}

\paragraph{Функция $g$ (замена расстояния при отсутствии соседа).}
\begin{equation}
  g(d,\, g_0,\, h) =
  \begin{cases}
    d,   & h = \text{истина (сосед есть)}, \\
    g_0, & h = \text{ложь}.
  \end{cases}
\label{eq:g}
\end{equation}
Здесь $d$ --- одно из $d_x^\pm$, $d_y^\pm$; $g_0$ --- подставляемая константа; $h$ --- флаг наличия соседа.

В реализации для оси $x$ всегда $g_0 = 0$; для оси $y$ --- $g_0 = w$.

\subsection{Сигналы ошибки}

Пусть $v_i = (v_x, v_y, v_z)^\top$ --- текущая скорость агента $i$ (\texttt{agent.dpose}). Определяются скаляры:
\begin{align}
  \sigma_x &= v_x - \xi\!\bigl(g(d_x^+,\, 0,\, h_x^+)\bigr) + \xi\!\bigl(g(d_x^-,\, 0,\, h_x^-)\bigr), \label{eq:sigma_x} \\
  \sigma_y &= v_y - \xi\!\bigl(g(d_y^+,\, w,\, h_y^+)\bigr) + \xi\!\bigl(g(d_y^-,\, w,\, h_y^-)\bigr), \label{eq:sigma_y}
\end{align}
где $h_x^+$, $h_x^-$, $h_y^+$, $h_y^-$ --- соответствующие флаги наличия соседа.

\paragraph{Интерпретация.}
Величины $\xi(g(\cdot))$ играют роль желаемых «вкладов'' в скорость, зависящих от относительных положений ближайших соседей справа/слева и сверху/снизу в системе координат агента. Разность с фактической скоростью образует ошибку $\sigma_x$, $\sigma_y$.

\subsection{Релейное управление}

Горизонтальная сила (ускорение при $m=1$):
\begin{equation}
  u_i = \bigl(u_x,\, u_y,\, 0\bigr)^\top,\qquad
  u_x = -u_{\max}\,\operatorname{sgn}(\sigma_x),\quad
  u_y = -u_{\max}\,\operatorname{sgn}(\sigma_y).
\label{eq:relay}
\end{equation}

\paragraph{Переменные.}
\begin{itemize}
  \item $\sigma_x$, $\sigma_y$ --- скалярные сигналы ошибки по осям.
  \item $u_x$, $u_y$ --- компоненты управления; по $Oz$ в этом сценарии управление нулевое, высота задаётся начальными условиями и не стабилизируется законом \eqref{eq:relay}.
\end{itemize}

Таким образом, каждая ось независимо подчиняется \emph{bang-bang} закону с насыщением по модулю $u_{\max}$.

%---------------------------------------------------------------------
\section{Динамика агента (интегрирование)}
\label{sec:dynamics}

После присвоения $F_i = u_i$ (\texttt{apply\_force}) выполняется шаг (\texttt{step}):

\paragraph{Переменные.}
\begin{itemize}
  \item $a_i = F_i / m_i$ --- ускорение (\texttt{ddpose}).
  \item Обновление скорости (явная схема Эйлера для $a$, затем обновление $v$):
  \begin{equation}
    v_i(t+\Delta t) = v_i(t) + a_i(t)\,\Delta t.
  \end{equation}
  \item Обновление положения (схема с учётом среднего ускорения в квадратичном члене):
  \begin{equation}
    p_i(t+\Delta t) = p_i(t) + v_i(t)\,\Delta t + \tfrac{1}{2} a_i(t)\,(\Delta t)^2.
  \end{equation}
\end{itemize}
Внешняя сила \texttt{external\_force} в вызове по умолчанию нулевая.

%---------------------------------------------------------------------
\section{Метрики для постобработки логов}

При анализе записанных траекторий строятся графики расстояния до оси коридора и «интервала'' вдоль оси.

\subsection{Геометрия коридора в плоскости $Oxy$}

Пусть $a_L = (x_L, y_L, z_L)^\top$, $a_R = (x_R, y_R, z_R)^\top$ --- положения якорей на текущем шаге. Вектор вдоль «оси коридора'' в плоскости:
\begin{equation}
  \mathbf{a} = \bigl(x_L - x_R,\; y_L - y_R\bigr)^\top \in \mathbb{R}^2.
\end{equation}

\subsection{Расстояние от агента до прямой через якоря (\texttt{plot\_graph\_axis\_err})}

Положение агента в плоскости: $\mathbf{p} = (x_i, y_i)^\top$. Вектор от правого якоря к агенту:
\begin{equation}
  \mathbf{b} = (x_R - x_i,\; y_R - y_i)^\top.
\end{equation}
Проекция $\mathbf{b}$ на направление $\mathbf{a}$ (со знаком, вдоль оси):
\begin{equation}
  p_{\parallel} = \frac{\mathbf{a}^\top \mathbf{b}}{\|\mathbf{a}\|}.
\end{equation}
Расстояние от точки $\mathbf{p}$ до прямой, проходящей через якорь $R$ в направлении $\mathbf{a}$:
\begin{equation}
  d_\perp = \sqrt{\|\mathbf{b}\|^2 - p_{\parallel}^2},
\end{equation}
что соответствует формуле длины перпендикуляра при не вырожденном $\mathbf{a}$.

\paragraph{Переменные.}
\begin{itemize}
  \item $\|\mathbf{a}\|$ --- длина вектора между якорями в плоскости.
  \item $\|\mathbf{b}\|$ --- расстояние от правого якоря до агента в плоскости.
  \item $d_\perp$ --- искомое расстояние до оси коридора; строится график $d_\perp(t)$ для каждого агента.
\end{itemize}

\subsection{«Интервал'' вдоль оси (\texttt{plot\_graph\_intervals})}

В коде используется величина
\begin{equation}
  d_{\mathrm{int}} = \frac{y_i}{\cos\varphi},\qquad
  \cos\varphi = \frac{a_y}{\|\mathbf{a}\|} = \frac{y_L - y_R}{\|\mathbf{a}\|},
\end{equation}
при $\cos\varphi \neq 0$. Это связано с параметризацией положения вдоль наклонной оси при определённой ориентации коридора; при движении якорей параллельно $Ox$ ($y_L \approx y_R$) знаменатель близок к нулю, и величина становится численно неустойчивой --- это ограничение конкретной формулы визуализации, а не закона управления.

%---------------------------------------------------------------------
\section{Логирование и визуализация}

\paragraph{Запись состояния агента в лог (\texttt{n\_params} = 12 полей).}
В каждой строке для каждого агента сохраняются: $(x,y,z)$, $(\dot x,\dot y,\dot z)$, $(u_x,u_y,u_z)$, \texttt{peer\_state}, $R_{\mathrm{vis}}$, признак \texttt{is\_alive}.

\paragraph{GIF (\texttt{create\_gif}).}
Строится анимация: траектории якорей, точки агентов, круги радиуса $R_{\mathrm{vis}}$ (заливка зоны восприятия).

\paragraph{Переменные визуализации.}
\begin{itemize}
  \item \texttt{frame\_step}, \texttt{start\_frame}, \texttt{end\_frame}, \texttt{n\_frames} --- прореживание кадров и ограничение длины анимации.
  \item \texttt{s\_point}, \texttt{s\_anchor} --- размеры маркеров на графике (в пикселях matplotlib \texttt{scatter}).
\end{itemize}

%---------------------------------------------------------------------
\section{Краткое резюме алгоритма (псевдокод)}

\begin{verbatim}
для n = 0 .. T-1:
    записать логи
    для каждого якоря k: t_k += dt; a_k <- f_k(t_k)
    для каждого живого агента i:
        P_i <- относительные векторы к видимым агентам и якорям
        вычислить (d_x^+, d_y^+, d_x^-, d_y^-) и флаги
        sigma_x, sigma_y по формулам из раздела «Закон управления»
        u_i <- (-u_max sgn(sigma_x), -u_max sgn(sigma_y), 0)
        интегрировать v_i, p_i на dt
\end{verbatim}

%---------------------------------------------------------------------
\section{Замечания о модели}

\begin{itemize}
  \item Управление \eqref{eq:relay} разрывное; в связке с $\tanh$ в \eqref{eq:sigma_x}--\eqref{eq:sigma_y} получается нелинейная обратная связь по относительным смещениям соседей.
  \item Якоря не имеют динамики второго порядка в этой симуляции: их положение задаётся кинематически.
  \item Высота $z$ не регулируется законом управления в \texttt{control\_force}; сценарий по сути плоский с «заглушкой'' по $z$.
\end{itemize}

\end{document}
